\documentclass[10pt,a4paper]{amsart}
\usepackage[margin=19mm]{geometry}
\usepackage{amsmath,amssymb,mathtools}
\usepackage[scr=rsfs,cal=euler]{mathalfa}
\usepackage{xfrac}
\usepackage[unicode,hidelinks,bookmarks=false]{hyperref}
\newcommand{\ie}{i.e.\ }
\newcommand{\cf}{cf.\ }
\newcommand{\resp}{respectively\ }
\newcommand{\Subsection}{\subsection}
\providecommand{\nobreakdash}{\nobreak\hbox{-}}
\setlength{\emergencystretch}{2em}
\multlinegap0pt
\usepackage{amssymb}
\usepackage{algorithm}\usepackage{algpseudocode}
\usepackage{dsfont}
\usepackage{float}
\newtheorem{proposition}{Proposition}
\title{Data-Driven Stochastic VRP: Integration of Forecast Duration into Optimization for Utility Workforce Management}
\author{Matteo Garbelli}
\date{Source: arXiv:2601.07514v1 (CC BY 4.0)}
\begin{document}
\maketitle

\noindent\emph{Source and licence.} Data-Driven Stochastic VRP: Integration of Forecast Duration into Optimization for Utility Workforce Management. Version arXiv:2601.07514v1 (CC BY 4.0), distributed by arXiv under the Creative Commons Attribution 4.0 International licence, http://creativecommons.org/licenses/by/4.0/, as recorded in the arXiv licence metadata for this exact version and shown on the abstract page https://arxiv.org/abs/2601.07514v1. Full licence notice: \texttt{licenses/arxiv-2601.07514-CC-BY-4.0.txt}.\par\smallskip

\noindent\textit{Editorial scope.} Counted unit: the article's proposition prop3.1 (source0.tex lines 269--278). The article's complete original proof of it is delivered with it (source0.tex lines 280--282, one \texttt{proof} environment of 3 lines). No mathematics is written, repaired or re-derived: the statement and the proof are the article's own lines, in the article's own notation, and no retained line is substituted or re-flowed.  No further source-local context is needed: every cross-reference of every retained line resolves inside the two delivered spans.  Nothing was omitted from any retained span, so the package carries no omission record.  No retained line contains a citation, so the wrapper carries no bibliography and no bibliography entry.  No retained line contains a figure environment, an image-inclusion command or drawing code, so no image is delivered and the package declares no asset dependency.  Editorial formatting only, in addition: an AMS \texttt{amsart} preamble that restates the article's own declarations and macros used by the retained lines, namely the environments \texttt{proposition} on separate counters, together with the standard packages those lines need.  The retained environments print in this document as Proposition 1 in order of appearance, whereas the article prints the same items with the numbering of its own full document.  The complete native source of the article as distributed in the arXiv e-print for this exact version is bundled unchanged as \texttt{sources/192507/source0.tex}.
\begin{proposition}[Route-level chance feasibility via sub-Gaussian buffer]
\label{prop3.1}
Let $R_k$ be the sequence of customers on vehicle $k$'s route and suppose $S_i=\mu_i+\varepsilon_i$ with independent sub-Gaussian $\varepsilon_i$ having proxy variances $\sigma_i^2$. If
\[
\sum_{i\in R_k}\mu_i + \sum_{(i,j)\in R_k} t_{ij}
+ \sqrt{2\log(1/\alpha_k)\sum_{i\in R_k}\sigma_i^2}
\ \le\ H_k,
\]
then $\mathbb{P}\{T_{n+1,k}-T_{0k} \le H_k\}\ge 1-\alpha_k$.
\end{proposition}

\begin{proof}
    By sub-Gaussian concentration, the sum of independent sub-Gaussian random variables $\sum_{i\in R_k}\varepsilon_i$ satisfies the tail bound above. The route duration is $\sum_{i\in R_k}(\mu_i+\varepsilon_i) + \sum_{(i,j)\in R_k} t_{ij}$. The deterministic part is $\sum_{i\in R_k}\mu_i + \sum_{(i,j)\in R_k} t_{ij}$. Adding the buffer $\Delta_{\alpha_k}(R_k)$ ensures that with probability at least $1-\alpha_k$, the stochastic deviations $\sum_{i\in R_k}\varepsilon_i$ do not exceed the buffer, hence the total route duration stays within $H_k$. \hfill$\square$
\end{proof}

\end{document}
